% Source: pca_overview.tex
\begin{notechapter}
% ---- chapter-local definitions ----
\tikzset{%
    myarrow/.style = {-Stealth, shorten >=5pt}
}
\forgetcs\mypoint\newcommand{\mypoint}[2]{\tikz[remember picture]{\node[inner sep=0, anchor=base](#1){$#2$};}}
\forgetcs\framedtext\newcommand{\framedtext}[1]{%
\par%
\noindent\fbox{%
    \parbox{\dimexpr\linewidth-2\fboxsep-2\fboxrule}{#1}%
}%
}
\forgetcs\horzbar\newcommand*{\horzbar}{\rule[.5ex]{5.5ex}{0.5pt}}
\forgetcs\vertbar\newcommand*{\vertbar}{\rule[-1ex]{0.5pt}{5.5ex}}
\guardnewlength{\Lnote}
\forgetcs\notte\newcommand{\notte}[1]
     {\addtolength{\leftmargini}{4em}
        \settowidth{\Lnote}{\textbf{Note:~}}
        \begin{quote}
            \rule{\dimexpr\textwidth-2\leftmargini}{1pt}\\
                        \mbox{}\hspace{-\Lnote}\textbf{Note:~}%
                                            #1\\[-0.5ex] 
            \rule{\dimexpr\textwidth-2\leftmargini}{1pt}
        \end{quote}
        \addtolength{\leftmargini}{-4em}}
\forgetcs\Tr\newcommand{\Tr}{\operatorname{Tr}}
% -----------------------------------
\chapter{Principal Component Analysis: Setting and Inertia}\label{ch:pca-overview}

\section{Setting}

Let's say that $X$ is p-dimensional random variable, e.g. $X = \begin{bmatrix*}X_1\\X_2\\ \vdots \\X_p \end{bmatrix*}$.\\[10pt]
\\[6pt]
We collect n samples and form the matrix $\textbf{X} \in \mathbb{R}^{n\times p}$:
\begin{align*}
    \textbf{X} &=
    \begin{bmatrix*}
        \tikzmarknode[circle]{left1}{x_{{1}{1}}}& \; x_{{1}{2}}&  \cdots \cdots  \quad &
        \tikzmarknode[circle]{right1}{x_{{1}{p}}} \mypoint{here1}{}\\
        x_{{2}{1}}& \; x_{{2}{2}}&  \cdots \cdots  \quad &x_{{2}{p}}\\
        \vdots& \; \vdots& \cdots \cdots  \quad &\vdots \\
        x_{{n}{1}}& \; x_{{n}{2}}&  \cdots \cdots  \quad &x_{{n}{p}}
    \end{bmatrix*}
\end{align*}
%
\begin{tikzpicture}[remember picture, overlay]
    \node[right=30pt of here1](textofhere1){One observation of $X = \begin{bmatrix*}X_1\\X_2\\ \vdots \\X_p \end{bmatrix*}$};
    \draw[myarrow] (textofhere1) -- (here1);
\end{tikzpicture}
%
\AddToShipoutPictureBG*{%
\begin{tikzpicture}[overlay,remember picture]
  \draw let  \p1=($(left1.north)-(left1.center)$),
  \p2=($(right1.north)-(right1.center)$), \n1={2*max(\y1,\y2)-2pt} in
  [opacity=.3,line width=\n1,line cap=round,
  shorten >=-\y2/3,shorten <=-\y1/3] (left1.center) -- (right1.center);
\end{tikzpicture}}
%
We can see $\textbf{X}$ as a collection of n individuals (observations) in $\mathbb{R}^{p}$ 
or we can see it as a collection of p variables in $\mathbb{R}^{n}$.
%
\begin{align*}
    \textbf{X}
    = 
        \begin{bmatrix*}
            {x}_{{1}{1}}& \cdots &\tikzmarknode[circle]{top3}{{x}_{{1}{k}}}&  \cdots  \quad & {x}_{{1}{p}} \\
            {x}_{{2}{1}}& \cdots &{x}_{{2}{k}}&  \cdots \quad &{x}_{{2}{p}}\\
            \vdots &  \cdots  &\vdots &\vdots \\
            \mypoint{}{\tikzmarknode[circle]{}{{x}_{{n}{1}}}}& \mypoint{here3}{\tikzmarknode[circle]{bottom3}{\cdots}}&  \mypoint{here3}{\tikzmarknode[circle]{bottom3}{x_{{n}{k}}}}&  \mypoint{}{\tikzmarknode[circle]{}{\cdots}}  \quad & \mypoint{}{\tikzmarknode[circle]{}{{x}_{{n}{p}}}}
        \end{bmatrix*}
\end{align*}
\begin{tikzpicture}[remember picture, overlay]
    \node[below=15pt of here3](textofhere3){sample values for the random variable ${X_k}$};
    \draw[myarrow] (textofhere3) -- (here3);
\end{tikzpicture}
%
\AddToShipoutPictureBG*{%
    \begin{tikzpicture}[overlay,remember picture]
      \draw let  \p1=($(top3.north)-(top3.center)$),
      \p2=($(bottom3.north)-(bottom3.center)$), \n1={2*max(\y1,\y2)-2pt} in
      [opacity=.3,line width=\n1,line cap=round,
      shorten >=-\y2/3,shorten <=-\y1/3] (top3.center) -- (bottom3.center);
    \end{tikzpicture}}
%
\subsection{Example}
%
\columnratio{0.3}
\begin{paracol}{2}

    \begin{figure}[h!]
        \includegraphics[width=100px]{X.png}
    \end{figure}
    
    \switchcolumn
    
    \begin{figure}[h!]
        \includegraphics[width=350px]{X_plots.png}
    \end{figure}
\end{paracol}
%
%\newpage
%
\subsection{Metrics}
We saw that we have two kind of vectors, the individuals (or observations),
$\mathbf{e}_{i} = \big[x_{{i}{1}} \, x_{{i}{2}} \cdots x_{{i}{p}} \big]^{T} \in \mathbb{R}^p \quad i \in \{1,\cdots,n\}$,\\[6pt]
and the variables, 
$\mathbf{x}_{k} = \big[x_{{1}{k}} \, x_{{2}{k}} \cdots x_{{n}{k}} \big]^{T} \in \mathbb{R}^n \quad k \in \{1,\cdots,p\}$.

\subsubsection{Metric in $\mathbb{R}^n$ - The individuals weight metric}
Let's consider two variables (i.e. two column-vectors of $\mathbf{X}$) $\mathbf{x}_{k} = \big[x_{{1}{k}} \, x_{{2}{k}} \cdots x_{{n}{k}} \big]^{T}$
and $\mathbf{x}_{l} = \big[x_{{1}{l}} \, x_{{2}{l}} \cdots x_{{n}{l}} \big]^{T}\quad k, l \in \{1,\cdots,p\}$.\\[6pt]
Each $x_{i}{.}$ of each vector represents a measurement of each individual.\\[5pt]
For example, $x_{{2}{k}}$ is the measurement of variable $k$ for individual $2$.\\[5pt]
%
If we consider a positive definite matrix $\mathbf{W}$, we can define an inner product in $\mathbb{R}^n$ as 
$\langle \mathbf{x}_{k}, \mathbf{x}_{l} \rangle_{\textbf{W}} = \mathbf{x}_{k}^{T}\textbf{W}\mathbf{x}_{l}$.\\[5pt]
There are many possibilities for $\mathbf{W}$, but one is particularly interesting in the case of centered data.\\[5pt]
Indeed if we take for $\mathbf{W}$ the weight matrix of individuals, that is \\[5pt]
$\textbf{W} = diag \big( \big[ w_1, \dots, w_i,\dots, w_n \big] \big) = 
\begin{bmatrix*}
    w_1 & 0    & \cdots  & \cdots  & 0 & \cdots & 0      \\
            0 & w_2 & 0  & \cdots  & \cdots & \cdots & 0      \\
        \vdots & \ddots & \ddots  & \ddots & \ddots & \ddots & \vdots \\
            0 & \cdots  & \ddots  & w_i & \ddots & \cdots & 0      \\
        \vdots & \ddots  & \ddots  & \ddots  & \ddots & \ddots & \vdots \\
        \vdots & \ddots  & \ddots  & \ddots  & \ddots & \ddots & 0      \\
            0 & \cdots    & \cdots  & \cdots  & \cdots & 0      & w_n
\end{bmatrix*}$, with $1 = \displaystyle\sum\limits_{i=1}^{n}w_i$. \\[10pt]
%
Then, for centered data,
$\langle \mathbf{x}_{k}, \mathbf{x}_{l} \rangle_{\textbf{W}} = \mathbf{x}_{k}^{T}\textbf{W}\mathbf{x}_{l} = s_{{k}{l}}$, 
the covariance of the two variables.\\[6pt]
%
The weight matrix $\mathbf{W}$ is the one we use for computing statistics on individuals.\\[5pt]
In the usual case where all individuals have the same weight $\mathbf{W} = \dfrac{1}{n}\mathbf{I}$.\\[5pt]
%
An example of statistics on individuals is the "mean" individual, that is, the gravity center $\mathbf{g}$ of individuals.\\[5pt]
%
$\mathbf{g}$ is given by the arithmetic mean $\displaystyle\sum\limits_{i=1}^{n}w_i\mathbf{e}_{i} = 
\begin{bmatrix*}\overline{x_{1}}\\\overline{x_{2}}\\ \vdots \\\overline{x_{p}} \end{bmatrix*}
$ where $\overline{x_{k}} = \displaystyle\sum\limits_{i=1}^{n}w_{i}x_{ik} = 
\big[ w_1 \dots w_i \dots w_n \big] \cdot \begin{bmatrix*}{x_{{1}{k}}}\\{x_{{2}{k}}}\\ \vdots \\{x_{{n}{k}}} \end{bmatrix*}
$.\\[6pt]
%
Which can be written in matrix notation $\mathbf{g} = \mathbf{X}^{T}\mathbf{W}\mathbf{1}_{n}$. 
Where $\mathbf{1}_{n} = \begin{bmatrix*}1\\1\\ \vdots \\1 \end{bmatrix*}$.
%
%
\subsubsection{Metric in $\mathbb{R}^p$}
In $\mathbb{R}^p$ we have $n$ individuals.\\[5pt]
$\mathbf{e}_{i} = \big[x_{{i}{1}} \, x_{{i}{2}} \cdots x_{{i}{p}} \big]^{T} \in \mathbb{R}^p \quad i \in \{1,\cdots,n\}$.\\[6pt]
%
We can compute the distance between individuals using an inner product.\\[5pt]
If $\textbf{M}$ is a positive definite matrix, then we have the scalar product of $\mathbf{e}_{i}$ and $\mathbf{e}_{j}$:
$\langle \mathbf{e}_{i}, \mathbf{e}_{j} \rangle_{\textbf{M}} = \mathbf{e}_{i}^{T}\textbf{M}\mathbf{e}_{j}$.\\[6pt]
Then the distance between two individuals $\mathbf{e}_{i}$ and $\mathbf{e}_{j}$ is given by 
$d_{\textbf{M}}^2(\mathbf{e}_{i}, \mathbf{e}_{j}) = 
\langle \mathbf{e}_{i} - \mathbf{e}_{j}, \mathbf{e}_{i} - \mathbf{e}_{j} \rangle_{\textbf{M}} = 
(\mathbf{e}_{i} - \mathbf{e}_{j})^{T}\textbf{M}(\mathbf{e}_{i} - \mathbf{e}_{j})$.
\\[5pt]
A usual choice for $\textbf{M}$ is the matrix of the variance inverses.\\[10pt]
$\textbf{M} = diag \big( \big[\dfrac{1}{s_{1}^2} \dots, \dfrac{1}{s_{k}^2} \dots, \dfrac{1}{s_{p}^2} \big] \big) = 
\begin{bmatrix*}
    \quad \dfrac{1}{s_{1}^2} & 0                                 & \cdots  & \cdots                                  & 0      & \cdots & 0      \\
                                        0 & \dfrac{1}{s_{2}^2} & 0       & \cdots                                  & \cdots & \cdots & 0      \\
                                    \vdots & \ddots                                  & \ddots  & \ddots                                  & \ddots & \ddots & \vdots \\
                                        0 & \cdots                                  & \ddots  & \dfrac{1}{s_{k}^2} & \ddots & \cdots & 0      \\
                                    \vdots & \ddots                                  & \ddots  & \ddots                                  & \ddots & \ddots & \vdots \\
                                    \vdots & \ddots                                  & \ddots  & \ddots                                  & \ddots & \ddots & 0      \\
                                        0 & \cdots                                  & \cdots  & \cdots                                  & \cdots & 0      & \dfrac{1}{s_{p}^2}
\end{bmatrix*}$.\\[6pt]
%
With $s_{k}^2 = \displaystyle\sum\limits_{i=1}^{n}w_{i}(x_{ik} - \overline{x_{k}})^2$,  
where $\overline{x_{k}} = \displaystyle\sum\limits_{i=1}^{n}w_{i}x_{ik}$.\\[6pt]
%
If $\mathbf{W} = \dfrac{1}{n}\mathbf{I}$ this becomes 
$s_{k}^2 = \dfrac{1}{n}\displaystyle\sum\limits_{i=1}^{n}(x_{ik} - \overline{x_{k}})^2$,  
and $\overline{x_{k}} = \dfrac{1}{n}\displaystyle\sum\limits_{i=1}^{n}x_{ik}$.
\\[10pt]
Then, the measurement units are suppressed and each variable has the same importance in the computation of distance.
%
\\
\notte{Since $\textbf{M}$ is positive definite it can be written as $\textbf{M} = \textbf{N}^{T}\textbf{N}$.\\[6pt]
Thus, $\langle \mathbf{e}_{i}, \mathbf{e}_{j} \rangle_{\textbf{M}} = \mathbf{e}_{i}^{T}\textbf{M}\mathbf{e}_{j} =
\mathbf{e}_{i}^{T}\textbf{N}^{T}\textbf{N}\mathbf{e}_{j} = \big(\textbf{N}\mathbf{e}_{i}\big)^{T}\big(\textbf{N}\mathbf{e}_{j}\big) =
\langle \textbf{N}\mathbf{e}_{i}, \textbf{N}\mathbf{e}_{j} \rangle_{\textbf{I}_{p}}$.\\[6pt]
Which is equivalent as using the identity matrix $I_{p}$ as metric on the transformed data $\textbf{X}\textbf{N}^{T}$.\\[6pt]
Thus, if $\textbf{M} = diag \big( \big[\dfrac{1}{s_{1}^2} \dots, \dfrac{1}{s_{k}^2} \dots, \dfrac{1}{s_{p}^2} \big] \big)$
then $\textbf{N} = diag \big( \big[\dfrac{1}{s_{1}} \dots, \dfrac{1}{s_{k}} \dots, \dfrac{1}{s_{p}} \big] \big)$.\\[6pt]
%
Then, the transformed data $\textbf{X}\textbf{N}^{T}$ is equal to the matrix $\textbf{X}$ where each variable is divided by its standard deviation.
}
%
\section{Study of the individuals}
We have n individuals $\mathbf{e}_{i}, i \in \{1,\cdots,n\}$, each of which is a row vector of $\mathbb{R}^{p}$.\\
\begin{align*}
    \textbf{X} &=
    \begin{bmatrix*}
        x_{{1}{1}}& \; x_{{1}{2}}&  \cdots \cdots  \quad &x_{{1}{p}}\\
        \vdots& \; \vdots& \cdots \cdots  \quad &\vdots \\
        x_{{i}{1}}& \; x_{{i}{2}}&  \cdots \cdots  \quad &x_{{i}{p}}\\
        \vdots& \; \vdots& \cdots \cdots  \quad &\vdots \\
        x_{{n}{1}}& \; x_{{n}{2}}&  \cdots \cdots  \quad &x_{{n}{p}}
    \end{bmatrix*}
    =
    \begin{bmatrix*}
        \horzbar &{\mathbf{e}_{1}}^T& \horzbar\\
        &\vdots&\\
        \horzbar &{\mathbf{e}_{i}}^T& \horzbar\\
        &\vdots&\\
        \horzbar &{\mathbf{e}_{n}}^T& \horzbar
    \end{bmatrix*}
\end{align*}
For individuals we want to express their similarity, that is, how close are they for each variables.
\\[5pt]
For two individuals $\mathbf{e}_{i}$ and $\mathbf{e}_{j}$ we have
$d_{\textbf{M}}^2(\mathbf{e}_{i}, \mathbf{e}_{j}) = 
\langle \mathbf{e}_{i} - \mathbf{e}_{j}, \mathbf{e}_{i} - \mathbf{e}_{j} \rangle_{\textbf{M}} = 
(\mathbf{e}_{i} - \mathbf{e}_{j})^{T}\textbf{M}(\mathbf{e}_{i} - \mathbf{e}_{j})$.
\\[10pt]
\begin{align*}
d_{\textbf{M}}^2(\mathbf{e}_{i}, \mathbf{e}_{j}) &= 
\begin{bmatrix*}
    x_{{i}{1}} - x_{{j}{1}} \quad x_{{i}{2}} - x_{{j}{2}} \quad  \cdots  \quad x_{{i}{p}} & - &x_{{j}{p}}
\end{bmatrix*}
\begin{bmatrix*}
    m_{{1}{1}}& \; m_{{1}{2}}&  \cdots \cdots  \quad &m_{{1}{p}}\\
    \vdots& \; \vdots& \cdots \cdots  \quad &\vdots \\
    m_{{l}{1}}& \; m_{{l}{2}}&  \cdots \cdots  \quad &m_{{l}{p}}\\
    \vdots& \; \vdots& \cdots \cdots  \quad &\vdots \\
    m_{{p}{1}}& \; m_{{p}{2}}&  \cdots \cdots  \quad &m_{{p}{p}}
\end{bmatrix*}
\begin{bmatrix*}
    x_{{i}{1}} & - & x_{{j}{1}} \\
    x_{{i}{2}} & - & x_{{j}{2}} \\
            & \vdots & \\
    x_{{i}{p}} & - &x_{{j}{p}}
\end{bmatrix*}\\
&=
\displaystyle\sum\limits_{l=1}^{p}\displaystyle\sum\limits_{k=1}^{p}m_{{l}{k}}(x_{{i}{l}} - x_{{j}{l}})(x_{{i}{k}} - x_{{j}{k}})
\end{align*}
\notte{Since $\textbf{M}$ defines an inner product it is symmetric (positive definite matrices are symmetric by definition), therefore $m_{kl} = m_{lk}$.}
%
For $\textbf{M} = diag \big( \big[\dfrac{1}{s_{1}^2} \dots, \dfrac{1}{s_{k}^2} \dots, \dfrac{1}{s_{p}^2} \big] \big)$ 
we get $d_{\textbf{M}}^2(\mathbf{e}_{i}, \mathbf{e}_{j}) = 
\displaystyle\sum\limits_{k=1}^{p}\dfrac{1}{{s}_{k}^2}(x_{{i}{k}} - x_{{j}{k}})^2$.
%
\subsection{Inertia}
To characterize the cloud of individuals we use inertia, that is the weighted sum of square distances between individuals.\\[5pt]
$\mathcal{I}^{*} = \dfrac{1}{2} 
\displaystyle\sum\limits_{i=1}^{n}\displaystyle\sum\limits_{j=1}^{n}{w_i}{w_j}d_{\textbf{M}}^2(\mathbf{e}_{i}, \mathbf{e}_{j})$.\\[5pt]
If all individuals are given the same weight, we get
$\mathcal{I}^{*} = \dfrac{1}{2n^2}
\displaystyle\sum\limits_{i=1}^{n}\displaystyle\sum\limits_{j=1}^{n}d_{\textbf{M}}^2(\mathbf{e}_{i}, \mathbf{e}_{j})$.\\[5pt]
\begin{prop}
$\mathcal{I}^{*} = \mathcal{I}_{\mathbf{g}} = 
\displaystyle\sum\limits_{i=1}^{n} {w_i} \big| \big| \mathbf{e}_{i} - \mathbf{g} \big| \big|_{\textbf{M}}^{2}$.
\end{prop}
\begin{proof}
\begin{align*}
    \mathcal{I}^{*} &= \dfrac{1}{2} 
    \displaystyle\sum\limits_{i=1}^{n}\displaystyle\sum\limits_{j=1}^{n}{w_i}{w_j}d_{\textbf{M}}^2(\mathbf{e}_{i}, \mathbf{e}_{j})\\
    &= \dfrac{1}{2} 
    \displaystyle\sum\limits_{i=1}^{n}\displaystyle\sum\limits_{j=1}^{n}{w_i}{w_j}
     \big| \big| {\textbf{e}_i - \textbf{g} + \textbf{g} - \textbf{e}_j\big| \big|}_{\mathbf{M}}^2\\
        &= \dfrac{1}{2} \bigg(
        \displaystyle\sum\limits_{i=1}^{n}\displaystyle\sum\limits_{j=1}^{n}{w_i}{w_j}
         \big| \big| {\textbf{e}_i - \textbf{g} \big| \big|}_{\mathbf{M}}^2 + 
        \displaystyle\sum\limits_{i=1}^{n}\displaystyle\sum\limits_{j=1}^{n}{w_i}{w_j}
         \big| \big| {\textbf{e}_j - \textbf{g} \big| \big|}_{\mathbf{M}}^2 - 2
        \displaystyle\sum\limits_{i=1}^{n}\displaystyle\sum\limits_{j=1}^{n}{w_i}{w_j}
         \langle \textbf{e}_i - \textbf{g}, \textbf{e}_j - \textbf{g} \rangle_{\textbf{M}} \bigg)\\
        &= \dfrac{1}{2} \bigg(
        \displaystyle\sum\limits_{j=1}^{n} {w_j} \displaystyle\sum\limits_{i=1}^{n}{w_i}
         \big| \big| {\textbf{e}_i - \textbf{g} \big| \big|}_{\mathbf{M}}^2 + 
        \displaystyle\sum\limits_{i=1}^{n} {w_i} \displaystyle\sum\limits_{j=1}^{n}{w_j}
         \big| \big| {\textbf{e}_j - \textbf{g} \big| \big|}_{\mathbf{M}}^2 - 2
            \Big \langle \displaystyle\sum_{i=1}^{n} w_{i} \mathbf{e}_{i} - \mathbf{g}\displaystyle\sum_{i=1}^{n} w_{i}, 
                \displaystyle\sum_{j=1}^{n} w_{j} \mathbf{e}_{j} - \mathbf{g}\displaystyle\sum_{j=1}^{n} w_{j}
            \Big \rangle_{\textbf{M}} \bigg)\\
        &= \dfrac{1}{2} \bigg( 
        \displaystyle\sum\limits_{i=1}^{n}{w_i}
         \big| \big| {\textbf{e}_i - \textbf{g} \big| \big|}_{\mathbf{M}}^2 +
        \displaystyle\sum\limits_{j=1}^{n}{w_j}
         \big| \big| {\textbf{e}_j - \textbf{g} \big| \big|}_{\mathbf{M}}^2 - 2
         \big \langle \mathbf{g} - \mathbf{g}, \mathbf{g} - \mathbf{g} \big \rangle_{\textbf{M}} \bigg)\\
        &= \dfrac{1}{2} \big(
        \mathcal{I}_{g} + \mathcal{I}_{g} \big) = \mathcal{I}_{g}
\end{align*}
\end{proof}
%
If $w_i = \dfrac{1}{n} \forall i$ then we obtain
$\mathcal{I}_{\mathbf{g}} = \dfrac{1}{n}
\displaystyle\sum\limits_{i=1}^{n} \big| \big| \mathbf{e}_{i} - \mathbf{g} \big| \big|_{\textbf{M}}^{2}$,
which is the mean of square distances of individuals to their gravity center.
\\[6pt]
For centered data, $\mathbf{g} = \mathbf{0}$ and $\mathcal{I} = \dfrac{1}{n}
\displaystyle\sum\limits_{i=1}^{n} \big| \big| \mathbf{e}_{i} \big| \big|_{\textbf{M}}^{2}$.\\[10pt]
We can show that $\mathcal{I}_{\mathbf{g}}$ is equal to $\Tr\big(\textbf{M} \textbf{V} \big)$.\\[10pt]
First, let's gather all $\mathbf{e}_{i} - \mathbf{g} \; (i \in \{1,\cdots,n\})$ in 
$\textbf{X}_{c}$ (the $n \times p$ matrix of centered data).\\[10pt]
$\textbf{X}_{c} = \textbf{X} - \mathbf{1}_{n}\mathbf{g}^{T} = \textbf{X} - \mathbf{1}_{n}\mathbf{1}_{n}^{T}\mathbf{W}\mathbf{X}
=
\begin{bmatrix*}
    \horzbar &{\mathbf{e}_{1}}^T - {\mathbf{g}}^T & \horzbar\\
    &\vdots&\\
    \horzbar &{\mathbf{e}_{i}}^T - {\mathbf{g}}^T & \horzbar\\
    &\vdots&\\
    \horzbar &{\mathbf{e}_{n}}^T - {\mathbf{g}}^T & \horzbar
\end{bmatrix*}
=
\begin{bmatrix*}
    \horzbar &[\textbf{e}_{1} - \textbf{g}]^T& \horzbar\\
    &\vdots&\\
    \horzbar &[\textbf{e}_{i} - \textbf{g}]^T& \horzbar\\
    &\vdots&\\
    \horzbar &[\textbf{e}_{n} - \textbf{g}]^T& \horzbar
\end{bmatrix*}
=
\begin{bmatrix*}
    \horzbar &\dot{\textbf{e}_{1}}^T& \horzbar\\
    &\vdots&\\
    \horzbar &\dot{\textbf{e}_{i}}^T& \horzbar\\
    &\vdots&\\
    \horzbar &\dot{\textbf{e}_{n}}^T& \horzbar
\end{bmatrix*} .$
\\[5pt]
Where $\mathbf{1}_{n} = \begin{bmatrix*}1\\1\\ \vdots \\1 \end{bmatrix*}$, $\mathbf{W}^{T}=\mathbf{W}$ 
and $\dot{\textbf{e}_{i}} = \textbf{e}_{i} - \textbf{g}$.\\[5pt]
Computing the covariance matrix of $\textbf{X}_{c}$ we get  $\textbf{V} = \textbf{X}_{c}^{T}\mathbf{W}\textbf{X}_{c} = 
\mathbf{X}^{T}\mathbf{W}\mathbf{X} - {\mathbf{g}}{\mathbf{g}}^T$\\[5pt]
%
\begin{align*}
\textbf{V} &= \begin{bmatrix*}
    \vertbar & \cdots & \vertbar & \cdots & \vertbar\\
    \dot{\textbf{e}_{1}} & \cdots & \dot{\textbf{e}_{i}} & \cdots & \dot{\textbf{e}_{n}}\\
    \vertbar & \cdots & \vertbar & \cdots & \vertbar
\end{bmatrix*}
\begin{bmatrix*}
    w_1 & 0    & \cdots  & \cdots  & 0  \\
    0 & \ddots & \ddots  & \ddots & \vdots \\
    \vdots & \ddots  &  w_i & \ddots & \vdots \\   
    \vdots & \ddots  & \ddots  & \ddots  & 0   \\
     0 & \cdots  & \cdots & 0      & w_n
\end{bmatrix*}
\begin{bmatrix*}
    \horzbar &{\dot{\textbf{e}_{1}}}^T & \horzbar\\
    &\vdots&\\
    \horzbar &{\dot{\textbf{e}_{i}}}^T & \horzbar\\
    &\vdots&\\
    \horzbar &{\dot{\textbf{e}_{n}}}^T & \horzbar
\end{bmatrix*}\\
&=
\begin{bmatrix*}
    \vertbar & \cdots & \vertbar & \cdots & \vertbar\\
    \dot{\textbf{e}_{1}} & \cdots & \dot{\textbf{e}_{i}} & \cdots & \dot{\textbf{e}_{n}}\\
    \vertbar & \cdots & \vertbar & \cdots & \vertbar
\end{bmatrix*}
\begin{bmatrix*}
    \horzbar & w_1 \dot{\textbf{e}_{1}}^T & \horzbar\\
    &\vdots&\\
    \horzbar & w_i \dot{\textbf{e}_{i}}^T & \horzbar\\
    &\vdots&\\
    \horzbar & w_n \dot{\textbf{e}_{n}}^T & \horzbar
\end{bmatrix*}
\end{align*}
\\[8px]
%
Now let's show that the inertia is equal to the trace of $\textbf{M}\textbf{V}$.\\[6pt]
%
Recall that since $\textbf{M}$ is definite positive, it can be written as $\textbf{M}=\textbf{N}^T \textbf{N}$.\\[10pt]
%
Thus, $w_{i} d_{\textbf{M}}^2(\textbf{e}_i,\textbf{g}) = 
w_{i} [\textbf{N}(\textbf{e}_i - \textbf{g})]^T [\textbf{N} (\textbf{e}_i - \textbf{g})] = w_i \textbf{y}_i^T \textbf{y}_i$. 
Where $\textbf{y}_i = \textbf{N}(\textbf{e}_i - \textbf{g}) = \textbf{N} \dot{\textbf{e}_i}$.\\[10pt]
%
%
And, $\textbf{M} \textbf{V} = \textbf{N}^T \textbf{N} \textbf{X}_c^T \textbf{W} \textbf{X}_c$, hence 
$\Tr \big(\textbf{M} \textbf{V} \big) = \Tr \big(\textbf{N}^T \textbf{N} \textbf{X}_c^T \textbf{W} \textbf{X}_c \big)$.\\[10pt]
%
Since $\Tr \big(\textbf{A} \textbf{B} \big) = \Tr \big(\textbf{B} \textbf{A} \big)$ we have 
$\Tr \big(\textbf{N}^T \textbf{N} \textbf{X}_c^T \textbf{W} \textbf{X}_c \big) = 
\Tr \big(\textbf{N} \textbf{X}_c^T \textbf{W} \textbf{X}_c \textbf{N}^T \big)$.\\[8pt]
%
Computing $\textbf{N}\textbf{X}_c^T$, we get
%
\begin{align*}
    \textbf{N}\textbf{X}_c^T =
    \begin{bmatrix*}
        \vertbar & \cdots & \vertbar & \cdots & \vertbar\\
        \textbf{N}\dot{e}_{1} & \cdots & \textbf{N}\dot{e}_{j} & \cdots & \textbf{N}\dot{e}_{n}\\
        \vertbar & \cdots & \vertbar & \cdots & \vertbar
    \end{bmatrix*}
    =
    \begin{bmatrix*}
        \vertbar & \cdots & \vertbar & \cdots & \vertbar\\
        {\textbf{y}}_{1} & \cdots & {\textbf{y}}_{j} & \cdots & {\textbf{y}}_{n}\\
        \vertbar & \cdots & \vertbar & \cdots & \vertbar
    \end{bmatrix*}
    = \textbf{Y}
\end{align*}
%
Thus, 
\begin{align*}
    \Tr\big(\textbf{M} \textbf{V} \big) = \Tr \big(\textbf{Y} \textbf{W} \textbf{Y}^T \big) &=
    \Tr \Bigg( \begin{bmatrix*}
        \vertbar & \cdots & \vertbar & \cdots & \vertbar\\
        {\textbf{y}}_{1} & \cdots & {\textbf{y}}_{j} & \cdots & {\textbf{y}}_{n}\\
        \vertbar & \cdots & \vertbar & \cdots & \vertbar
    \end{bmatrix*} 
    \textbf{W} 
    \begin{bmatrix*}
        \horzbar &\textbf{y}_{1}^T& \horzbar\\
        &\vdots&\\
        \horzbar &\textbf{y}_{i}^T& \horzbar\\
        &\vdots&\\
        \horzbar &\textbf{y}_{n}^T& \horzbar
    \end{bmatrix*}\Bigg) \\
    &= \Tr \Bigg( \begin{bmatrix*}
        \vertbar & \cdots & \vertbar & \cdots & \vertbar\\
        {\textbf{y}}_{1} & \cdots & {\textbf{y}}_{j} & \cdots & {\textbf{y}}_{n}\\
        \vertbar & \cdots & \vertbar & \cdots & \vertbar
    \end{bmatrix*} 
    \begin{bmatrix*}
        w_1(\horzbar &\textbf{y}_{1}^T& \horzbar)\\
        &\vdots&\\
        w_i(\horzbar &\textbf{y}_{i}^T& \horzbar)\\
        &\vdots&\\
        w_n(\horzbar &\textbf{y}_{n}^T& \horzbar)
    \end{bmatrix*}\Bigg)\\
    &= \Tr \Bigg(\sum_{k=1}^{n} 
    \begin{bmatrix*}
        \vertbar\\
        {\textbf{y}}_{k} \\
        \vertbar
    \end{bmatrix*}
    \begin{bmatrix*}
        w_k(\horzbar &\textbf{y}_{k}^T& \horzbar)
    \end{bmatrix*}\Bigg) \\
    &= \sum_{k=1}^{n} \Tr \Bigg(\begin{bmatrix*}
        \vertbar\\
        {\textbf{y}}_{k} \\
        \vertbar
    \end{bmatrix*}
    \begin{bmatrix*}
        w_k(\horzbar &\textbf{y}_{k}^T& \horzbar)
    \end{bmatrix*}\Bigg)  \\
    &= \sum_{k=1}^{n} \Tr \Bigg(
    \begin{bmatrix*}
        w_k(\horzbar &\textbf{y}_{k}^T& \horzbar)
    \end{bmatrix*}  
    \begin{bmatrix*}
        \vertbar\\
        {\textbf{y}}_{k} \\
        \vertbar
    \end{bmatrix*} \Bigg) \\
    &= \displaystyle\sum_{k=1}^{n} w_k \textbf{y}_k^T \textbf{y}_k = \displaystyle\sum_{k=1}^{n} w_{k} d_{\textbf{M}}^2(\textbf{e}_k,\textbf{g}) 
\end{align*}
\\[10pt]
We want to project the individuals in an affine subspace of $\mathbb{R}^{p}$ such that the cloud of projected individuals
is "as close as possible" to the original cloud of individuals.\\[6pt]
%
Let's define what we mean by "close".\\[6pt]
First, we introduce some notations. \\[6pt]
We will denote $\mathbb{V}\big[\textbf{u}_1, \cdots, \textbf{u}_{q}\big]$ the
subspace of $\mathbb{R}^{p}$ generated by the orthonormal vectors $\big\{\textbf{u}_1, \cdots, \textbf{u}_{q}\big\}$ of $\mathbb{R}^{p}$ ($q \leq p$).\\[6pt]
Then an affine space of origin $\textbf{a}$ parallel to $\mathbb{V}\big[\textbf{u}_1, \cdots, \textbf{u}_{q}\big]$ is noted 
$E_{\textbf{a}} = \textbf{a} + \mathbb{V}\big[\textbf{u}_1, \cdots, \textbf{u}_{q}\big]$.\\[6pt]
The $\textbf{M}$-orthogonal projection of an individual $\textbf{e}_i$ on $E_{\textbf{a}}$ is denoted $\hat{\textbf{e}}_i$.\\[6pt]
By "as close as possible", we mean that 
$\mathcal{J}_{E_{\textbf{a}}} =  \displaystyle\sum_{i=1}^{n} w_{i}d_{\mathbf{M}}^2(\textbf{e}_i, \hat{\textbf{e}}_i)$
must be as small as possible.\\[6pt]
\begin{prop}
    $\mathcal{J}_{E_{\textbf{g}}} \leq \mathcal{J}_{E_{\textbf{a}}}$ for every $\textbf{a} \in \mathbb{R}^p$, with equality if and only if $\textbf{g} \in E_{\textbf{a}}$ (i.e.\ $E_{\textbf{a}} = E_{\textbf{g}}$).
    In words: among all affine subspaces with a given direction $\mathbb{V}\big[\textbf{u}_1, \cdots, \textbf{u}_{q}\big]$, the best one goes through the gravity center.
\end{prop}
\begin{proof}
    Let's denote $\mathcal{H}$ the set of all affine spaces of direction $\mathbb{V}\big[\textbf{u}_1, \cdots, \textbf{u}_{q}\big]$.\\[6pt]
    That is $\mathcal{H} = \big\{ \textbf{a} + \mathbb{V}\big[\textbf{u}_1, \cdots, \textbf{u}_{q}\big], \forall \textbf{a} \in \mathbb{R}^{p} \big\}$.\\[6pt]
    Let $H \in \mathcal{H}$, and  $H_{\textbf{g}} = \textbf{g} + \mathbb{V}\big[\textbf{u}_1, \cdots, \textbf{u}_{q}\big]$ 
    (subspace parallel to $H$ going through $\textbf{g}$).\\[6pt]
    Let's denote $\hat{\textbf{g}}$ the projection of $\textbf{g}$ on $H$.\\[6pt]
    Since $\textbf{g} - \hat{\textbf{g}}$ is $\textbf{M}$-orthogonal to the direction of $H$, the projection on $H_{\textbf{g}}$ is the projection on $H$ translated by $(\textbf{g} - \hat{\textbf{g}})$.\\[5pt]
    Thus,
    \begin{align*}
        \mathcal{J}_{H_{\textbf{g}}} &=
        \displaystyle\sum_{i=1}^{n} w_{i}d_{\mathbf{M}}^2(\textbf{e}_i, \hat{\textbf{e}}_i + (\textbf{g} - \hat{\textbf{g}}))\\
        &=
        \displaystyle\sum_{i=1}^{n} w_{i} \big| \big| {\textbf{e}_i - (\hat{\textbf{e}}_i + (\textbf{g} - \hat{\textbf{g}})) \big| \big|}_{\mathbf{M}}^2\\
        &=
        \displaystyle\sum_{i=1}^{n} w_{i} \big| \big| {\textbf{e}_i - \hat{\textbf{e}}_i \big| \big|}_{\mathbf{M}}^2 + 
        \displaystyle\sum_{i=1}^{n} w_{i} \big| \big| {\textbf{g} - \hat{\textbf{g}} \big| \big|}_{\mathbf{M}}^2 - 2
        \displaystyle\sum_{i=1}^{n} w_{i} \langle \mathbf{e}_{i} - \hat{\mathbf{e}_{i}}, \mathbf{g} - \hat{\mathbf{g}} \rangle_{\textbf{M}}\\
        &=
        \displaystyle\sum_{i=1}^{n} w_{i} \big| \big| {\textbf{e}_i - \hat{\textbf{e}}_i \big| \big|}_{\mathbf{M}}^2 + 
        \displaystyle\sum_{i=1}^{n} w_{i} \big| \big| {\textbf{g} - \hat{\textbf{g}} \big| \big|}_{\mathbf{M}}^2 - 2
        \langle \displaystyle\sum_{i=1}^{n} w_{i} \big( \mathbf{e}_{i} - \hat{\mathbf{e}_{i}} \big), \mathbf{g} - \hat{\mathbf{g}} \rangle_{\textbf{M}}\\
        &=
        \displaystyle\sum_{i=1}^{n} w_{i} \big| \big| {\textbf{e}_i - \hat{\textbf{e}}_i \big| \big|}_{\mathbf{M}}^2 + 
        \displaystyle\sum_{i=1}^{n} w_{i} \big| \big| {\textbf{g} - \hat{\textbf{g}} \big| \big|}_{\mathbf{M}}^2 - 2
        \langle \displaystyle\sum_{i=1}^{n} w_{i} \mathbf{e}_{i} - \displaystyle\sum_{i=1}^{n} w_{i}\hat{\mathbf{e}_{i}}, \mathbf{g} - \hat{\mathbf{g}} \rangle_{\textbf{M}}\\
        &=
        \displaystyle\sum_{i=1}^{n} w_{i} \big| \big| {\textbf{e}_i - \hat{\textbf{e}}_i \big| \big|}_{\mathbf{M}}^2 + 
        \displaystyle\sum_{i=1}^{n} w_{i} \big| \big| {\textbf{g} - \hat{\textbf{g}} \big| \big|}_{\mathbf{M}}^2 - 2
        \langle \mathbf{g} - \hat{\mathbf{g}}, \mathbf{g} - \hat{\mathbf{g}} \rangle_{\textbf{M}}\\
        &=
        \bigg(\displaystyle\sum_{i=1}^{n} w_{i} \big| \big| {\textbf{e}_i - \hat{\textbf{e}}_i \big| \big|}_{\mathbf{M}}^2 \bigg) + 
        \big| \big| {\textbf{g} - \hat{\textbf{g}} \big| \big|}_{\mathbf{M}}^2 - 2
        \langle \mathbf{g} - \hat{\mathbf{g}}, \mathbf{g} - \hat{\mathbf{g}} \rangle_{\textbf{M}}\\
        &=
        \bigg(\displaystyle\sum_{i=1}^{n} w_{i} \big| \big| {\textbf{e}_i - \hat{\textbf{e}}_i \big| \big|}_{\mathbf{M}}^2 \bigg) -
        \big| \big| {\textbf{g} - \hat{\textbf{g}} \big| \big|}_{\mathbf{M}}^2\\
        &=
        \mathcal{J}_{H} - \big| \big| {\textbf{g} - \hat{\textbf{g}} \big| \big|}_{\mathbf{M}}^2
   \end{align*}
\end{proof}
%
Therefore, the optimal affine subspace $H_{opt}$ (in the sense of minimizing $\mathcal{J}_{H}$) has to be sought
among the affine subspaces that go through $\textbf{g}$.
%
\notte{In the demonstration above, we used the fact that $\displaystyle\sum_{i=1}^{n} w_{i}\hat{\mathbf{e}_{i}} = \hat{\mathbf{g}}$.\\[5pt]
That is, the gravity center of the projected observations is the projection of the gravity center of the observations.\\[5pt]
This is true since any projection $P$ is an affine application (i.e. $P$ preserves the barycenters).\\[5pt]
This is obvious, since $\hat{\mathbf{g}} = P\big(\displaystyle\sum_{i=1}^{n} w_{i}\mathbf{e}_{i}\big) = 
\displaystyle\sum_{i=1}^{n} w_{i} P(\mathbf{e}_{i}) =
\displaystyle\sum_{i=1}^{n} w_{i}\hat{\mathbf{e}_{i}}$.\\}
%
\begin{prop}
   The affine subspace $H_{opt}$ maximizes the inertia of the projected cloud on $H$:\\[6pt]
   \begin{align*}
        \mathcal{I}_{H_{opt}} = \underset{H \in \mathcal{H}_{\textbf{g}}}{\max}\ \mathcal{I}_{H}
   \end{align*}
   where $\mathcal{H}_{\textbf{g}} \subset \mathcal{H}$ is the set of all affine spaces going through $\textbf{g}$, and 
   $ \mathcal{I}_{H} = 
   \displaystyle\sum_{i=1}^{n} w_{i} \big| \big| {\hat{\textbf{e}}_i - \hat{\textbf{g}} \big| \big|}_{\mathbf{M}}^2$
   is the inertia of the cloud projected $\textbf{M}$-orthogonally on $H$.
\end{prop}
\begin{proof}
Let $H \in \mathcal{H}_{\textbf{g}}$, then $\hat{\textbf{g}} = \textbf{g}$ and using the Pythagorean theorem we have:
\begin{align*}
    \mathcal{J}_{H} = \displaystyle\sum_{i=1}^{n} w_{i} {\big| \big| \textbf{e}_i - \hat{\textbf{e}}_i \big| \big|}_{\mathbf{M}}^2
    = 
    \displaystyle\sum_{i=1}^{n} w_{i} 
    \big( {\big| \big| \textbf{e}_i - \textbf{g} \big| \big|}_{\mathbf{M}}^2 - 
    {\big| \big| \textbf{g} - \hat{\textbf{e}}_i \big| \big|}_{\mathbf{M}}^2 \big)
    =
    \mathcal{I}_{g} - \mathcal{I}_{H}
\end{align*}
Therefore, $\mathcal{I}_{H} = \mathcal{I}_{g} - \mathcal{J}_{H}$ is maximum for $\mathcal{J}_{H}$ minimum. That is, for $H = H_{opt}$.
\end{proof}
%
\begin{prop}
    For $H \in \mathcal{H}_{\textbf{g}}$, the weighted average of the square distances between points of the projected cloud,
    $\displaystyle\sum\limits_{i=1}^{n}\displaystyle\sum\limits_{j=1}^{n}{w_i}{w_j}d_{\textbf{M}}^2(\hat{\mathbf{e}}_{i}, \hat{\mathbf{e}}_{j})$,
    is maximal for $H = H_{opt}$.
\end{prop}
\begin{proof}
This is straightforward, since we have 
$\dfrac{1}{2} \displaystyle\sum\limits_{i=1}^{n}\displaystyle\sum\limits_{j=1}^{n}{w_i}{w_j}
d_{\textbf{M}}^2(\hat{\mathbf{e}}_{i}, \hat{\mathbf{e}}_{j}) 
= \displaystyle\sum_{i=1}^{n} w_{i} \big| \big| {\hat{\textbf{e}}_i - \hat{\textbf{g}} \big| \big|}_{\mathbf{M}}^2$\\
(see the proof that $\mathcal{I}^{*} = \mathcal{I}_{\mathbf{g}}$ above, applied to the projected cloud whose gravity center is $\hat{\mathbf{g}}$).
\end{proof}
%

To sum-up these last three propositions, the least deformation of the cloud of observation ($\mathcal{J}_{E_{a}}$) by projection 
on an affine space is obtained by maximizing the inertia of the projected cloud, which is equivalent to maximizing
the weighted average of the squared distances between points of the projected cloud. 

\subsection{Sequential resolution}
It remains to find the optimal subspace through $\mathbf{g}$. Since
$\mathcal{J}_{H} = \mathcal{I}_{\mathbf{g}} - \mathcal{I}_{H}$, we maximize the
inertia of the projected cloud.

\paragraph{One axis.}
Let $H = \mathbf{g} + \mathbb{V}[\mathbf{u}]$ with $\|\mathbf{u}\|_{\mathbf{M}}^2 = \mathbf{u}^T\mathbf{M}\mathbf{u} = 1$.
The $\mathbf{M}$-orthogonal projection of $\mathbf{e}_i$ on $H$ is
$\hat{\mathbf{e}}_i = \mathbf{g} + \langle \dot{\mathbf{e}}_i, \mathbf{u} \rangle_{\mathbf{M}}\, \mathbf{u}$, hence
\[
\mathcal{I}_H = \sum_{i=1}^n w_i \langle \dot{\mathbf{e}}_i, \mathbf{u}\rangle_{\mathbf{M}}^2
= \sum_{i=1}^n w_i\, \mathbf{u}^T \mathbf{M} \dot{\mathbf{e}}_i \dot{\mathbf{e}}_i^T \mathbf{M} \mathbf{u}
= \mathbf{u}^T \mathbf{M} \mathbf{V} \mathbf{M} \mathbf{u}.
\]
Maximizing $\mathbf{u}^T \mathbf{M}\mathbf{V}\mathbf{M}\mathbf{u}$ under the constraint
$\mathbf{u}^T\mathbf{M}\mathbf{u} = 1$ with the Lagrangian
$\mathcal{L}(\mathbf{u},\lambda) = \mathbf{u}^T \mathbf{M}\mathbf{V}\mathbf{M}\mathbf{u} - \lambda(\mathbf{u}^T\mathbf{M}\mathbf{u} - 1)$
gives $\mathbf{M}\mathbf{V}\mathbf{M}\mathbf{u} = \lambda \mathbf{M}\mathbf{u}$, i.e.\ (since $\mathbf{M}$ is invertible)
\[
\mathbf{V}\mathbf{M}\,\mathbf{u} = \lambda\, \mathbf{u},
\qquad\text{and then}\qquad \mathcal{I}_H = \mathbf{u}^T\mathbf{M}(\lambda\mathbf{u}) = \lambda .
\]
The first principal axis $\mathbf{u}_1$ is therefore the eigenvector of $\mathbf{V}\mathbf{M}$ associated
with its largest eigenvalue $\lambda_1$, and $\lambda_1$ is the inertia it explains.

\paragraph{The spectrum of $\mathbf{V}\mathbf{M}$.}
$\mathbf{V}\mathbf{M}$ is $\mathbf{M}$-symmetric: $\mathbf{M}(\mathbf{V}\mathbf{M}) = (\mathbf{V}\mathbf{M})^T\mathbf{M}$.
Writing $\mathbf{M} = \mathbf{N}^T\mathbf{N}$, the matrix $\mathbf{N}\mathbf{V}\mathbf{N}^T$ is symmetric positive
semi-definite and $\mathbf{V}\mathbf{M}\mathbf{u} = \lambda\mathbf{u} \iff \mathbf{N}\mathbf{V}\mathbf{N}^T(\mathbf{N}\mathbf{u}) = \lambda(\mathbf{N}\mathbf{u})$.
Hence $\mathbf{V}\mathbf{M}$ is diagonalizable, its eigenvalues
$\lambda_1 \ge \lambda_2 \ge \cdots \ge \lambda_p \ge 0$ are real and non-negative, and its
eigenvectors can be chosen $\mathbf{M}$-orthonormal.

\paragraph{Next axes.}
If $\mathbf{u}_1,\dots,\mathbf{u}_q$ are $\mathbf{M}$-orthonormal, the projection on
$\mathbf{g}+\mathbb{V}[\mathbf{u}_1,\dots,\mathbf{u}_q]$ is the sum of the projections on the axes, and by
Pythagoras $\mathcal{I}_H = \sum_{k=1}^q \mathbf{u}_k^T\mathbf{M}\mathbf{V}\mathbf{M}\mathbf{u}_k$.
Looking for $\mathbf{u}_2$ maximizing $\mathbf{u}^T\mathbf{M}\mathbf{V}\mathbf{M}\mathbf{u}$ under
$\mathbf{u}^T\mathbf{M}\mathbf{u}=1$ and $\mathbf{u}^T\mathbf{M}\mathbf{u}_1 = 0$ gives the eigenvector associated
with $\lambda_2$, and so on. The sequential (greedy) solution is also the globally optimal
$q$-dimensional one: by Ky Fan's maximum principle applied to $\mathbf{N}\mathbf{V}\mathbf{N}^T$, the maximum
of $\mathcal{I}_H$ over all $q$-dimensional subspaces through $\mathbf{g}$ is $\lambda_1 + \dots + \lambda_q$,
attained by $\mathbb{V}[\mathbf{u}_1,\dots,\mathbf{u}_q]$.

\paragraph{Summary.}
\begin{itemize}
  \item The principal axes are the $\mathbf{M}$-orthonormal eigenvectors $\mathbf{u}_k$ of $\mathbf{V}\mathbf{M}$.
  \item The total inertia is $\mathcal{I}_{\mathbf{g}} = \Tr(\mathbf{V}\mathbf{M}) = \sum_{k=1}^p \lambda_k$, and
        axis $k$ explains the fraction $\lambda_k / \sum_j \lambda_j$ of it.
  \item The principal components (coordinates of the individuals on the axes) are
        $\mathbf{c}_k = \mathbf{X}_c \mathbf{M} \mathbf{u}_k \in \mathbb{R}^n$, with weighted variance
        $\mathbf{c}_k^T \mathbf{W} \mathbf{c}_k = \lambda_k$.
  \item With $\mathbf{M} = \mathbf{I}_p$ we recover the usual PCA on the covariance matrix; with
        $\mathbf{M} = \operatorname{diag}(1/s_1^2,\dots,1/s_p^2)$ we obtain the normed PCA (PCA on the correlation matrix).
\end{itemize}

\paragraph{Further reading.} G.~Saporta, \emph{Probabilit\'es, analyse des donn\'ees et statistique};
B.~Escofier and J.~Pag\`es, \emph{Analyses factorielles simples et multiples}.

\end{notechapter}
